% -*- Document-Type: LaTeX; Master-File: manual.TX -*- 

% \chapter{A short introduction to \elan}\label{short}

This chapter presents in a top-down approach the \elan\ main
features. If you are beginning to use \elan, you should
certainly have a look at this chapter first. Complementarily
Chapter~\ref{language} presents a bottom-up complete
description of the language.

For a gently introduction to \elan\ we recommend to
read~\cite{BorovanskyKKMR-WRLA98} or to access to 
\localisationWebElan{}.

%%%%%%%%%%%%%%%%%%%%%%%%%%%%%%%%%%%%%%%%%
\section{What could you do in \elan?}
%%%%%%%%%%%%%%%%%%%%%%%%%%%%%%%%%%%%%%%%%

Relying on the premiss that inferences rules can be quite conveniently
described by rewrite rules, we started in the early nineties to design
and implement a language in which inference systems can be represented
in a natural way, and executed reasonably efficiently.

This leads us quickly to formalise such a language using the rewriting
logic introduced by J.~Meseguer~\cite{MeseguerTCS92} and to see \elan\
as a logical framework where the frame logic is rewriting logic
extended with the fundamental notion of
strategies~\cite{VittekThese,KirchnerKV-MIT95}.  A rewrite theory and
an associated strategy is called a {\em computational system}.

In \elan, a logic can be expressed by specifying its syntax, its
inference rules and a description of how these rules should be
applied. The syntax of the logic can be described using mixfix
operators as in the OBJ language~\cite{obj3-88} or
SDF~\cite{HeeringHendriksKlintRekers89}.  The inference rules of the
logic are described by conditional rewrite rules with \texttt{where}
assignments allowing to introduce variables local to a rule.

From this description of the logic and a theory written in this logic,
we infer   automatically  a  computational   system  consisting of   a
rewriting theory plus strategies.

Since we wanted \elan{} to be modular, with a syntactic level
well-adapted to the user's needs, the language is designed in such a way
that three programming levels are possible.

\begin{itemize}
\item  First   the  design of  a    logic is  done  by  the  so-called
        \Def{super-user}, with the help of a powerful preprocessor that
        expands specific macros into \elan\ constructions.
\item Such a logic can be used by the (standard) \Def{user} in order
        to write a specification.
\item Finally,  the   \Def{end-user}  can  evaluate  queries  in   the
        specification, following the semantics described by the logic.
\end{itemize}
This corresponds to the general diagram given in Figure~\ref{elan-principle}.
The query is interactively given by the end-user at the \elan\ prompt
level.

\begin{figure}[hbt]
\begin{center}
\framebox{\parbox{11cm}{%\textwidth}{
$$\epsfig{figure=elan.1, width=10cm}$$
\caption{Elan: general execution schema \label{elan-principle}}
}}
\end{center}
\end{figure}

%A simple example, fully described later in this user manual, is the
%design of syntactic unification, where the tasks are divided as
%follows:
%\begin{enumerate}
%\item the super-user describes in a generic way the unification
%        inferences, i.e. a logic for unification,
%\item the user gives the specification of an algebra in which (s)he
%        wants to unify terms; in this case, this is quite simple since
%        it amounts to specify the function symbols of the 
%        considered term algebra,
%\item the end-user gives a unification problem.
%\end{enumerate}
%The diagram of Figure~\ref{elan-principle} thus instantiates  as
%described in Figure~\ref{elan-principle-unif}.

%\begin{figure}[hbt]
%\begin{center}
%\framebox{\parbox{11cm}{%\textwidth}{
%$$\epsfig{figure=elan.2, width=10cm}$$
%\caption{Elan: how to implement syntactic unification \label{elan-principle-unif}}
%}}
%\end{center}
%\end{figure}

%The description of the logic and of the specification are done in the
%\elan\ syntax described in the Chapter~\ref{language} which can be
%further extended by the super-user. These descriptions are
%assumed to be done in files with specific extensions:
%\begin{enumerate}
%\item{\texttt{.lgi}}  for the top level logic description, this file is written
%by the \elan\ super-user,
%\item{\texttt{.eln}}  for a module used in a logic description, this file is
%also written
%by the \elan\ super-user,
%\item{\texttt{.spc}}  for a specification, i.e. a program written in the
%defined logic by an \elan\ user. 
%\end{enumerate}

%The query is interactively given by the end-user at the \elan\ prompt
%level.

%\commentCK{Attention langage de commande}

%%%%%%%%%%%%%%%%%%%%%%%%%%%%%%%%%%%%%%%%%
\section{A very simple example}
%%%%%%%%%%%%%%%%%%%%%%%%%%%%%%%%%%%%%%%%%

Let us fully detail a complete example.
The next module illustrates what the programmer
has to write in order to describe the derivative operation on simple
polynomials. The first part contains:
\begin{itemize}
\item the module name: \texttt{poly1}
\item modules you want  to  import: \texttt{int} (the  integer library
        module)
\item sort declaration: \texttt{variable} and \texttt{poly}
\item operators definitions: constructors and functions declarations
\end{itemize}

The second part contains the computation definition.
Given a query, \elan\ repeatedly normalises the term using unlabelled
rules. 
This is done in order to perform functional evaluation and thus
it is recommended for the user to provide a confluent and terminating
unlabelled rewrite system to ensure termination and unicity of the
result. This normalisation process is built in the evaluation
mechanism and consists in a leftmost innermost normalisation. This
yields always a {\em single} result.

You can control the \elan\ construction of terms by giving
parsing annotations like associativity or priority to operators.
\insertExamples{poly1.eln}

The top level of the logic description given in
the following module describes the way to run the
system. The logic declaration just introduces the sorts of interest
(query and result) and defines \texttt{.eln} modules to be imported:
this is the ``top level''.
\insertExamples{poly1.lgi}

Once defined, we can use the two previous modules by calling the
interpreter as detailed in section~\ref{Run the interpreter}.

\vbox{
{\small\begin{verbatim}
% elan poly1.lgi
enter query term finished by the key word 'end': 
  deriv(X) end

[] start with term: deriv(X)
[] result term:       1

enter query term finished by the key word 'end':
  deriv(3*X*X + 2*X + 7) end

[] start with term: deriv(3*X*X+2*X+7)
[] result term:       0*X*X+3*1*X+X*1+0*X+2*1+0
\end{verbatim}}
}

%%%%%%%%%%%%%%%%%%%%%%%%%%%%%%%%%%%%%%%%%%
\section{A more generic example}
%%%%%%%%%%%%%%%%%%%%%%%%%%%%%%%%%%%%%%%%%%

The next example, introduces more \elan\ features.  It consists of the
specification of  elementary polynomials  built  over a finite set  of
variables, and the derivative  functions to compute the derivated form
with respect to a given variable.
%%%
Tasks are divided as follows:
\begin{enumerate}
\item the  super-user describes in a  generic  way the derivative  and
        factorisation   inferences,  i.e.    a  logic  for  polynomial
        differentiation,
\item the user gives the specification  of  an algebra in which  (s)he
        wants to  derivate polynomials;  in this  case, this is  quite
        simple  since it  amounts to   specify  the variables of   the
        considered polynomial algebra,
\item the end-user gives a differentiation problem.
\end{enumerate}

The diagram in Figure~\ref{elan-principle} thus instantiates  as
described in Figure~\ref{elan-principle-derivation}.

\begin{figure}[hbt]
\begin{center}
\framebox{\parbox{11cm}{%\textwidth}{
$$\epsfig{figure=elan.3, width=10cm}$$
\caption{Elan: how to implement polynomial differentiation 
  \label{elan-principle-derivation}}
}}
\end{center}
\end{figure}

The description of the logic and of the specification is done in the
\elan\ syntax described in the Chapter~\ref{language} and it can be
further extended by the super-user. These descriptions are
assumed to be done in files with specific extensions:
\begin{enumerate}
\item{\texttt{.lgi}} for the top level logic description, this file is
        written by the \elan\ super-user,
\item{\texttt{.eln}}  for a module used  in a  logic description, this
        file is also written by the \elan\ super-user,
\item{\texttt{.spc}} for a  specification, i.e.  a program written  in
        the defined logic by an \elan\ user.
\end{enumerate}
%%%

The \texttt{.eln} module of our example is the following:
\insertExamples{poly2.eln}

Then, the logic declaration
describes the way to parse the specification file that contains a
finite list of variables.
\insertExamples{poly2.lgi}

To instantiate the \textit{Logic} and build
the \textit{Computational System}, you have to give a specification.
An example of such specification is:
\insertExamples{someVariables.spc}

It contains a list of variables that is read and transformed into a
rewrite rule: \texttt{Vars => X.Y.Z.nil}

The \texttt{FOR EACH} preprocessor construction used in the
\texttt{poly2.eln} module performs a textual
replacement that extracts identifiers \texttt{X}, \texttt{Y},
\texttt{Z} from the list \texttt{X.Y.Z.nil} and declare them of sort   
\texttt{variable}.

You can notice that operators \texttt{*} and \texttt{+} are now
declared as associative and commutative, which is helpful to
implement some simplification rules.
One rule contains a conditional part.
        %Dollar signs `\$' are only used to solve parsing ambiguities.

Now let us call \elan\ with the specification
\texttt{someVariables.spc}: 

\vbox{{\small\begin{verbatim}
% elan poly2 someVariables.spc
enter query term finished by the key word 'end':
  deriv(3*X*X + 2*X + 7 , X) end

[] start with term: deriv(3*X*X+2*X+1,X)
[] result term:       X+X*3+2
\end{verbatim}}}

The result has been simplified in a better way, but it is still difficult to
read due to the lack of parenthesis.

%%%%%%%%%%%%%%%%%%%%%%%%%%%%%%%%%%%%%%%%%%%%%
\section{An extended example}
%%%%%%%%%%%%%%%%%%%%%%%%%%%%%%%%%%%%%%%%%%%%%

This last example shows how to play with labelled rules and
strategies.


Unlabelled rules are applied repeatedly, at any position, to normalize
the term. A labelled rule is applied only when a strategy calls it. In
this case, the rule is applied at top position on the term. In the
next module, you can recognize labelled rules: names are given between
brackets.


You can define strategies in  the same way  you  define rules and  use
them  in local  affectation     parts introduced  by   the   key  word
\texttt{where}. When   a  rule is    applied, before constructing  the
right-hand-side,    conditional  parts  are   evaluated  and variables
involved in local affectation parts are instantiated  by the result of
the application of a strategy on a term.  
\insertExamples{poly3.eln}
\label{poly.example}

And now we can use it:
{\small
\begin{verbatim}
% elan poly3 someVariables.spc
enter query term finished by the key word 'end':
  deriv(3*X*X + 2*X + 7 , X) end

[] start with term: deriv(3*X*X+2*X+1,X)
[] result term:       ((X*6)+2)
[] end
\end{verbatim}
} 


%%% Local Variables: 
%%% mode: latex
%%% TeX-master: t
%%% TeX-master: t
%%% End: 
